\documentclass[11pt]{article}
\usepackage{amsmath, amssymb}
\usepackage{geometry}
\geometry{margin=1in}
\usepackage{fancyhdr}
\usepackage{lastpage}
\usepackage{parskip} % Better paragraph spacing
\usepackage{titlesec} % Better section formatting
\usepackage{graphicx}
\usepackage{float}
\usepackage[most]{tcolorbox}
\usepackage{lmodern}
\usepackage{caption}
\usepackage{hyperref}
\usepackage{listings}
\usepackage{xcolor}

\definecolor{codegreen}{rgb}{0,0.6,0}
\definecolor{codegray}{rgb}{0.5,0.5,0.5}
\definecolor{codepurple}{rgb}{0.58,0,0.82}
\definecolor{backcolour}{rgb}{0.95,0.95,0.92}

% Define code styling
\lstdefinestyle{mystyle}{
    backgroundcolor=\color{backcolour},
    commentstyle=\color{codegreen},
    keywordstyle=\color{magenta},
    numberstyle=\tiny\color{codegray},
    stringstyle=\color{codepurple},
    basicstyle=\ttfamily\footnotesize,
    breakatwhitespace=false,
    breaklines=true,
    captionpos=b,
    keepspaces=true,
    numbers=left,
    numbersep=5pt,
    showspaces=false,
    showstringspaces=false,
    showtabs=false,
    tabsize=2
}

\lstset{style=mystyle}

% Define commands for reusable content
\newcommand{\coursename}{CS-451 Distributed Algorithms}
\newcommand{\lecturenum}{01}
\newcommand{\lecturedate}{07/09/2026}
\newcommand{\professor}{Rachid Guerraoui}

\hypersetup{colorlinks=true,linkcolor=blue,urlcolor=blue}

% Configure section spacing
\titlespacing*{\section}{0pt}{2ex plus 1ex minus .2ex}{1.5ex plus .2ex}
\titlespacing*{\subsection}{0pt}{1.5ex plus 1ex minus .2ex}{1ex plus .2ex}
% Header / Footer
\pagestyle{fancy}
\fancyhf{}
\fancyhead[L]{\coursename}
\fancyhead[R]{Lecture \lecturenum}
\fancyfoot[L]{EPFL - Computer Science}
\fancyfoot[C]{Author: Jonathan Dutoit}
\fancyfoot[R]{Page \thepage\ of \pageref{LastPage}}

\begin{document}

\begin{center}
  {\large EPFL - Computer Science} \\[0.5em]
  {\Large \textbf{\coursename}} \\[0.5em]
  {\large Lecture \lecturenum\ --- \lecturedate} \\[0.5em]
  {\large Author: Jonathan Dutoit} \\[0.5em]
  {\large Professor: \professor} \\[0.5em]
\end{center}

\vspace{2cm}

\begin{center}
	{\LARGE \textbf{Overview}}
\end{center}
\vspace{0.5cm}

\begin{tcolorbox}[
    enhanced,
    frame hidden,
    borderline west={2pt}{0pt}{black},
    boxsep=5pt,
    left=8pt
]
\textbf{Definition.} A \textbf{Distributed system} is a set of processes (machines), which are Turing machines, that communicate through shared messages.  

\end{tcolorbox}
\textbf{Ex.:} Google. Hundreds of thousands of machines connected. One job can involve 2000 machines, and on average 10 machines go down per day. 

The application needs underlying services for distributed interaction. The network \textcolor{red}{is not enough}. Reliability guarantees (e.g., TCP) are only offered for communication among pairs of processes, i.e., \textit{one-to-one communication (client-server)}.

\subsection*{Underlying services}
\begin{enumerate}
	\item \textbf{\textit{processes}} (abstracting computers)
	\item \textbf{\textit{channels}} (abstracting networks)
	\item \textbf{\textit{failure detectors}} (abstracting time)
\end{enumerate}

\begin{tcolorbox}[
    enhanced,
    frame hidden,
    borderline west={2pt}{0pt}{black},
    boxsep=5pt,
    left=8pt
]
\textbf{Definition.} Processes. The distributed system is made of a finite set of \textit{\textbf{processes}}: each process models a \textbf{\textit{sequential}} program.
\[
	\text{Processes are denoted} \quad  p_1, ..., p_N \quad \text{or \textit{p}, \textit{q}, \textit{r}}
.\] 

Processes have unique \textbf{\textit{identities}} and know each other. Every pair of processes is connected by a \textbf{\textit{link}} through which the processes exchanges \textbf{\textit{messages}}. \\

Every process has a \textbf{\textit{piece}} of the distributed algorithm. Modules within the same process interact by exchanging events.\\

\textbf{\textit{Event based programming}}:
\begin{lstlisting}[language=Octave]
upon event < Event1, att1, att2, ...> do 
	// something
	trigger < Event2, att1, att2, ...>
\end{lstlisting}
\end{tcolorbox}

\begin{tcolorbox}[
    enhanced,
    frame hidden,
    borderline west={2pt}{0pt}{black},
    boxsep=5pt,
    left=8pt
]
\textbf{Definition.} \textbf{\textit{Safety}} is a property which states that nothing bad should happen. \textbf{\textit{Liveness}} is a property which states that something good should happen. \\

\textbf{Example}: \textit{Tell the truth}
\begin{itemize}
	\item Having to say something is \textit{liveness}
	\item Not lying is \textit{safety}
\end{itemize}
\end{tcolorbox}

\subsection*{Fair-loss links}
\begin{itemize}
	\item \textbf{\textit{FL1. Fair-loss}}. If a message is sent infinitely often by pi to pj, and neither pi or pj crashes, them m is delivered infinitely often by pj $\Rightarrow$ \textcolor{blue}{liveness}.
	\item \textit{\textbf{FL2. Finite duplication}}. If a message m is sent a finite number of times by pi to pj, m is delivered a finite number of times by pj. $\Rightarrow$ \textcolor{blue}{liveness}. 
	\item \textbf{\textit{FL3. No creation}}. No message is delivered unless it was sent. $\Rightarrow$ \textcolor{blue}{safety}.
\end{itemize}

\subsection*{Stubborn links}
\begin{itemize}
	\item \textbf{\textit{SL1. Stubborn delivery}}. If a process pi sends a message m to a correct process pj, and pi does not crash, then pj delivers m an infinite number of times.
	\item \textbf{\textit{SL2. No creation}}. No message is delivered unless it was sent. 
\end{itemize}

\textbf{Implements:} SturbbornLinks (sp2p). \\
\textbf{Uses:} FairLossLinks (flp2p).
\begin{lstlisting}[language=Octave]
upon event <sp2pSend, dest, m> do
	while (true) do
		trigger <flp2pSend, dest, m>; 
	upon event <flp2pDeliver, src, m> do
		trigger (sp2pDeliver, src, m>;
\end{lstlisting}

\textbf{Note:} by the underlying link (fair-loss links), we know that if we send a message an infinite number of times, then the message is delivered infinitely. For delivering, the message comes first from the underlying link (fair-loss), and then it's transferred to the upper link (stubborn). 

\subsection*{Reliable (Perfect) links}

\begin{itemize}
	\item \textbf{\textit{PL1. Validity}}. If pi and pj are correct, then every message sent by pi to pj is eventually delivered by pj. 
	\item \textbf{\textit{PL2. No duplication}}. No message is delivered (to a process) more than once.
	\item \textbf{\textit{PL3. No creation}}. No message is delivered unless it was sent. 
\end{itemize}
\textbf{Implements}: PerfectLinks (pp2p) \\
\textbf{Uses}: StubbornLinks (sp2p).
\begin{lstlisting}[language=Octave, mathescape=true]
upon event < Init> do delivered := $\varnothing$
upon event < pp2pSend, dest, m>; 
	trigger < sp2pSend, dest, m>; 
upon event < sp2pDeliver, src, m> do
	if m $\notin$ delivered then
		trigger < pp2pDeliver, src, m>;
		add m to delivered;
\end{lstlisting}

Roughly speaking, reliable links ensure that messages exchanged between correct processes are \textbf{\textit{not lost}}.
\end{document}
